%%
%% ICT Express (elsarticle, 3p twocolumn)
%%
\documentclass[final,3p,times,twocolumn]{elsarticle}

\usepackage{graphicx, amsmath, pifont, amssymb}
\usepackage{enumitem}
\usepackage{booktabs}
\usepackage{array}
\usepackage{tabularx}
\usepackage{float}
\usepackage{placeins}
\usepackage[utf8]{inputenc}
\usepackage[english]{babel}
\usepackage{hyperref}
\hypersetup{colorlinks=true, linkcolor=black, filecolor=black, urlcolor=cyan}
\usepackage{caption}
\captionsetup{justification=raggedright, singlelinecheck=false}
\captionsetup[table]{labelformat=simple, labelsep=newline}
\captionsetup[figure]{labelformat=simple, labelsep=period}

% --- ProVerif listings ---
\usepackage{listings}
\usepackage{xcolor}
\lstdefinelanguage{ProVerif}{%
  morekeywords={fun,reduc,equation,event,query,free,type,channel,%
    forall,let,in,if,then,else,new,out,process,phase,not,data,%
    private,bitstring,attacker,inj,true,false},
  sensitive=true,
  morecomment=[s]{(*}{*)}}
\lstset{%
  language=ProVerif,
  basicstyle=\scriptsize\ttfamily,
  keywordstyle=\bfseries\color{blue!70!black},
  commentstyle=\color{gray!80!black}\itshape,
  frame=single,
  rulecolor=\color{gray!50},
  framesep=2pt,
  xleftmargin=3pt,
  xrightmargin=3pt,
  framexleftmargin=3pt,
  breaklines=true,
  columns=flexible,
  keepspaces=true,
  showstringspaces=false,
  aboveskip=4pt,
  belowskip=2pt,
  captionpos=b}

\biboptions{sort&compress}

\graphicspath{{figs/}{plots/}{pqc_benchmark_data/plots/}}

\journal{ICT Express}

\begin{document}

\begin{frontmatter}

\title{Evaluating Post-Quantum Cryptography for Consumer Remote SIM
Provisioning: Benchmarks and Migration Outlook\tnotemark[t1]}

\tnotetext[t1]{This work evaluates the performance impact and deployment
feasibility of post-quantum cryptography within the GSMA SGP.22 Consumer
Remote SIM Provisioning protocol.}

\author[1]{Author Name 1\corref{cor1}}
\ead{author1@example.org}

\author[1]{Author Name 2}
\ead{author2@example.org}

\address[1]{Organization, Department, City, Postcode, Country}

\cortext[cor1]{Corresponding author}

\begin{abstract}
The GSMA SGP.22 Remote SIM Provisioning (RSP) standard secures eSIM profile
delivery with ECDSA, ECDH, and X.509 certificate chains — all broken by
Shor's algorithm.  NIST has standardised ML-KEM (FIPS~203) and ML-DSA
(FIPS~204) as replacements, yet their larger key and signature sizes raise
open questions on constrained eUICC hardware.  We present the first
systematic PQC evaluation within the Consumer RSP lifecycle, benchmarking
ML-KEM-768 and ML-DSA-44 across all 29 SGP.22 protocol steps and analysing
four incremental migration configurations.  Results quantify latency, message
size, and certificate chain overhead, and support a phased migration outlook
for Consumer RSP operators.
\end{abstract}

\begin{keyword}
benchmarking \sep eSIM \sep ML-DSA \sep ML-KEM \sep migration strategy \sep
post-quantum cryptography \sep remote SIM provisioning \sep SGP.22
\end{keyword}

\end{frontmatter}

% ===================================================================
%  SECTION 1 — INTRODUCTION
% ===================================================================
\section{Introduction}\label{sec:intro}

The GSMA SGP.22 specification~\citep{gsma_sgp22_v26} standardises Remote SIM
Provisioning (RSP) for consumer devices, enabling over-the-air eSIM profile
delivery without physical SIM handling.  Every provisioning transaction depends
on ECDSA mutual authentication and ECDH session key establishment — primitives
that Shor's algorithm breaks in polynomial time on a cryptographically relevant
quantum computer (CRQC)~\citep{shor1994algorithms}.  The harvest-now-decrypt-later
(HNDL) threat means the exposure window opens before any CRQC
exists~\citep{mosca2018cybersecurity}, since provisioned profiles carry long-lived
IMSI and authentication keys.

NIST standardised ML-KEM and ML-DSA in 2024~\citep{nist2024pqc} as the
primary replacements, but their larger keys and signatures inflate SGP.22
certificate chains transmitted over metered ES9+ channels to constrained eUICC
hardware (typically $<$500\,KB RAM).  Prior benchmarks cover PQC in
TLS~\citep{schwabe2020kemtls,sikeridis2020pqtls}, 5G NAS~\citep{hoque2025pqc5g},
and secure-element channels~\citep{bettale2025pqesim}, but none jointly address
the SGP.22 provisioning flow, eUICC memory constraints, and a staged migration
path. This paper closes that gap.

\noindent\textbf{Contributions:}
\begin{itemize}[leftmargin=*, noitemsep, topsep=2pt]
  \item First per-step benchmark of ML-KEM-768 and ML-DSA-44 across all
        29 SGP.22 protocol steps, quantifying computational latency, certificate
        chain inflation, and over-the-air bandwidth for four migration options.
  \item A phased migration outlook grounded in benchmark results, identifying
        each option's deployment prerequisites and hardware dependencies.
\end{itemize}

% ===================================================================
%  SECTION 2 — BACKGROUND
% ===================================================================
\section{Background}\label{sec:background}

\subsection{Consumer RSP Architecture and Protocol}\label{sec:bg-arch}

SGP.22 Consumer RSP involves four entities: the SM-DP+ server (profile
preparation and delivery), the eUICC (tamper-resistant secure element), the
LPA (host-device mediator), and the GSMA Certificate Issuer (CI) trust
anchor.  The LPA connects to the SM-DP+ over the HTTPS ES9+ interface;
the eUICC never directly terminates the TLS session.
A provisioning session proceeds in two cryptographic phases.
\textbf{Phase~1 (Mutual Authentication, steps 1--16):} a dual-challenge,
dual-signature exchange where both parties sign session-specific structures
(\texttt{serverSigned1}/\texttt{euiccSigned1}) against CI certificate chains.
\textbf{Phase~2 (Profile Download, steps 17--29):} the SM-DP+ signs profile
metadata, performs an ephemeral key agreement with the eUICC, and delivers the
AES-128-protected Bound Profile Package (BPP).
Figures~\ref{fig:rsp-auth}--\ref{fig:rsp-download} detail both phases across
the four migration options evaluated here.

\begin{figure*}[p]
\centering
\includegraphics[width=\textwidth, height=0.44\textheight, keepaspectratio]{figs/CONSUMER-RSP-auth.png}
\caption{Phase~1: Mutual authentication (steps 1--16).
Options~(a)(b) use ECDSA; (c)(d) use ML-DSA.
ES9+ uses classical TLS in (a) and PQ-TLS in (b)--(d).}
\label{fig:rsp-auth}
\vspace{1em}
\includegraphics[width=\textwidth, height=0.44\textheight, keepaspectratio]{figs/CONSUMER-RSP-download.png}
\caption{Phase~2: Profile download (steps 17--29).
Options~(a)(b)(c) use ECDH for key agreement; (d) replaces it with ML-KEM.
Options~(a)(b) sign with ECDSA; (c)(d) with ML-DSA.}
\label{fig:rsp-download}
\end{figure*}

\subsection{Classical Primitives and Quantum Threat}\label{sec:bg-classical}

SGP.22 mandates ECDSA (P-256/P-384) for authentication and ECDH (P-256) for
session key establishment via ECIES; the resulting shared secret is fed into
the X9.63 KDF~\citep{nist_sp800_56a} to derive AES-128 BSP keys~\citep{nist_sp800_38b}.
Table~\ref{tab:classical_primitives} lists these primitives.
Both ECDSA and ECDH derive security from the elliptic-curve discrete logarithm
problem, which Shor's algorithm breaks simultaneously~\citep{proos2003shor};
AES-128 and SHA-256 face only Grover's quadratic speedup and are retained
in all four migration options.

\begin{table}[!htb]
\caption{SGP.22 classical primitives. ``Shor?'' marks schemes broken by a CRQC.}
\label{tab:classical_primitives}
\centering
\footnotesize
\setlength{\tabcolsep}{3pt}
\renewcommand{\arraystretch}{1.05}
\begin{tabularx}{\columnwidth}{@{}l >{\raggedright\arraybackslash}X cc@{}}
\toprule
\textbf{Alg.} & \textbf{Role} & \textbf{bits} & \textbf{Shor?} \\
\midrule
ECDSA P-256/P-384 & Mutual authentication & 128/192 & Y \\
ECDH P-256 & Session key (ECIES) & 128 & Y \\
AES-128 & BPP encryption (BSP) & 128 & No \\
SHA-256 & KDF, certificate digests & 128 & No \\
\bottomrule
\end{tabularx}
\end{table}

\subsection{NIST PQC Standards: ML-KEM and ML-DSA}\label{sec:bg-pqc}

ML-KEM (FIPS~203, from CRYSTALS-Kyber) is an IND-CCA2-secure Key
Encapsulation Mechanism based on Module-LWE~\citep{nist_fips203}.
ML-DSA (FIPS~204, from CRYSTALS-Dilithium) is an EUF-CMA-secure signature
scheme based on Module-SIS/LWE~\citep{nist_fips204}.  We use ML-KEM-768
(NIST Level~3, matching the P-384 option in SGP.22) to replace ECDH and
ML-DSA-44 (NIST Level~2) to replace ECDSA; ML-DSA-44 minimises
per-certificate overhead while remaining well above the current 128-bit
classical baseline~\citep{demir2025pqc}.
Table~\ref{tab:pqc_sizes} shows the size impact central to our bandwidth
analysis.

\begin{table}[!htb]
\caption{Sizes (bytes): pk, signature or KEM ciphertext, secret key.}
\label{tab:pqc_sizes}
\centering
\footnotesize
\setlength{\tabcolsep}{4pt}
\renewcommand{\arraystretch}{1.05}
\begin{tabular}{@{}lrrr@{}}
\toprule
\textbf{Algorithm} & \textbf{pk} & \textbf{sig\,/\,ct} & \textbf{sk} \\
\midrule
ECDSA-P256  & 64  & 64   & 32 \\
ECDSA-P384  & 96  & 96   & 48 \\
ML-DSA-44   & 1{,}312 & 2{,}420 & 2{,}560 \\
\midrule
ECDH-P256   & 64  & 64   & 32 \\
ECDH-P384   & 96  & 96   & 48 \\
ML-KEM-768  & 1{,}184 & 1{,}088 & 2{,}400 \\
\bottomrule
\end{tabular}
\end{table}

% ===================================================================
%  SECTION 3 — RELATED WORK
% ===================================================================
\section{Related Work}\label{sec:related}

PQC overhead in TLS has been studied extensively:
\cite{schwabe2020kemtls} benchmarked KEMTLS, \cite{sikeridis2020pqtls}
measured certificate-size-dominated latency under varying network conditions,
and \cite{gonzalez2025pqtls} provided a reusable OQS-based evaluation
framework~\cite{oqs2024} — but none address application-layer protocols
above TLS.  \cite{hoque2025pqc5g} analysed PQC in 5G NAS, yet 5G
authentication runs on gigabyte-RAM application processors, not the
$<$500\,KB eUICC.  \cite{ahmed2024rsp} formally verified SGP.22 security
properties without any performance measurement.  Closest to ours,
\cite{bettale2025pqesim} benchmarked PQC variants of GlobalPlatform SCP11
on a Cortex-M3 secure element, but SCP11 governs application-management
channels, not the SGP.22 provisioning lifecycle with its GSMA CI hierarchy
and profile-binding mechanism.  \cite{demir2025pqc} supplied per-operation
timings without protocol-flow mapping.
Table~\ref{tab:related_work} summarises coverage; this paper is the first
to jointly address all five dimensions.

\begin{table}[!htb]
\caption{Coverage matrix.
  PQ\,=\,PQC eval.; RSP\,=\,Consumer RSP; HW\,=\,HW bench.;
  SE\,=\,secure-element class; Mig.\,=\,migration outlook;
  FV\,=\,formal verification.}
\label{tab:related_work}
\centering
\footnotesize
\setlength{\tabcolsep}{3pt}
\renewcommand{\arraystretch}{1.05}
\begin{tabular}{@{}lcccccc@{}}
\toprule
 & \textbf{PQ} & \textbf{RSP} & \textbf{HW} & \textbf{SE} & \textbf{Mig.} & \textbf{FV} \\
\midrule
Schwabe et al.~\cite{schwabe2020kemtls} & \checkmark & & & & & \\
Hoque et al.~\cite{hoque2025pqc5g} & \checkmark & & & & $\sim$ & \\
Bettale et al.~\cite{bettale2025pqesim} & \checkmark & & \checkmark & \checkmark & & \\
Demir et al.~\cite{demir2025pqc} & \checkmark & & \checkmark & $\sim$ & & \\
Ahmed et al.~\cite{ahmed2024rsp} & & \checkmark & & & & \checkmark \\
\textbf{This work} & \checkmark & \checkmark & \checkmark & \checkmark & \checkmark & \checkmark \\
\bottomrule
\end{tabular}
\end{table}

% ===================================================================
%  SECTION 4 — SYSTEM MODEL AND MIGRATION CONFIGURATIONS
% ===================================================================
\section{System Model}\label{sec:system}

We adopt a Dolev-Yao adversary~\citep{dolev1983security} with quantum
capability (HNDL model): the attacker can record all ES9+ traffic and
eventually run Shor's algorithm against ECDSA/ECDH; the CI root and eUICC
hardware are trusted; replay is mitigated by SGP.22 session-unique challenges.
SGP.22 v2.6~\citep{gsma_sgp22_v26} is the baseline; only certificate and
signature field sizes change — ASN.1/TLV framing is preserved.

\subsection{Four Migration Options}\label{sec:sys-configs}

Table~\ref{tab:configs} defines four configurations.
\textbf{(a)}~Fully classical baseline.
\textbf{(b)}~Server-side PQ-TLS (ML-KEM-768 handshake) with no eUICC change —
the LPA terminates TLS, so immediate HNDL protection on the transport channel
is deployable today.
\textbf{(c)}~ML-DSA-44 replaces ECDSA for all authentication and binding
signatures; CI roots re-issued; new eUICC silicon with ML-DSA, but not ML-KEM.
\textbf{(d)}~Full application-layer PQC: ML-DSA-44 authentication plus
ML-KEM-768 session key agreement; eUICC requires both primitives.

\begin{table}[!htb]
\caption{Migration options.  BPP uses AES-128-BSP unchanged in all options.}
\label{tab:configs}
\centering
\footnotesize
\setlength{\tabcolsep}{3pt}
\renewcommand{\arraystretch}{1.05}
\begin{tabularx}{\columnwidth}{@{} >{\raggedright\arraybackslash}p{0.34\columnwidth}
    *{4}{>{\centering\arraybackslash}X} @{}}
\toprule
 & \textbf{(a)} & \textbf{(b)} & \textbf{(c)} & \textbf{(d)} \\
\midrule
SM-DP+/eUICC certs & ECDSA & ECDSA & ML-DSA-44 & ML-DSA-44 \\
Session key agr. & ECDH & ECDH & ECDH & ML-KEM-768 \\
ES9+ transport & TLS~1.3 & PQ-TLS & PQ-TLS & PQ-TLS \\
CI roots & ECDSA & ECDSA & ML-DSA-44 & ML-DSA-44 \\
\bottomrule
\end{tabularx}
\end{table}

\subsection{Cryptographic Substitution Mapping}\label{sec:sys-mapping}

Table~\ref{tab:substitution} maps each changing SGP.22 sub-step to the
algorithm used per option.  Options~(a) and~(b) are identical at the
application layer.
\emph{Sign}$_E$/\emph{Vrfy}$_E$ = ECDSA sign/verify;
\emph{Sign}$_M$/\emph{Vrfy}$_M$ = ML-DSA;
\emph{KA}$_E$ = ECDH shared secret;
\emph{kg/enc/dec} = ML-KEM key-gen/encaps/decaps.

\begin{table}[!htb]
\caption{Per-step crypto substitution across options (a)--(d).}
\label{tab:substitution}
\centering
\footnotesize
\setlength{\tabcolsep}{2pt}
\renewcommand{\arraystretch}{1.05}
\begin{tabularx}{\columnwidth}{@{}c >{\raggedright\arraybackslash}X ccc@{}}
\toprule
\textbf{St.} & \textbf{Operation} & \textbf{(a)(b)} & \textbf{(c)} & \textbf{(d)} \\
\midrule
8e  & SM-DP+ signs \texttt{serverSigned1}       & Sign$_E$ & Sign$_M$ & Sign$_M$ \\
12c & eUICC verifies \texttt{serverSignature1}   & Vrfy$_E$ & Vrfy$_M$ & Vrfy$_M$ \\
13b & eUICC signs \texttt{euiccSigned1}          & Sign$_E$ & Sign$_M$ & Sign$_M$ \\
16d & SM-DP+ verifies \texttt{euiccSignature1}   & Vrfy$_E$ & Vrfy$_M$ & Vrfy$_M$ \\
17c & SM-DP+ signs \texttt{smdpSigned2}          & Sign$_E$ & Sign$_M$ & Sign$_M$ \\
20b & eUICC verifies \texttt{smdpSignature2}     & Vrfy$_E$ & Vrfy$_M$ & Vrfy$_M$ \\
21a & eUICC ephemeral key                        & ECDH & ECDH & ML-KEM kg \\
24b & SM-DP+ session KA                          & KA$_E$   & KA$_E$   & ML-KEM enc \\
24h & SM-DP+ signs \texttt{bind\_body}           & Sign$_E$ & Sign$_M$ & Sign$_M$ \\
27a & eUICC verifies \texttt{serverSignature3}   & Vrfy$_E$ & Vrfy$_M$ & Vrfy$_M$ \\
27b & eUICC session KA                           & KA$_E$   & KA$_E$   & ML-KEM dec \\
27f & eUICC signs installation result            & Sign$_E$ & Sign$_M$ & Sign$_M$ \\
\bottomrule
\end{tabularx}
\end{table}

% ===================================================================
%  SECTION 5 — BENCHMARKING METHODOLOGY
% ===================================================================
\section{Benchmarking Methodology}\label{sec:methodology}

We measure per-operation latency, certificate chain sizes, and end-to-end
provisioning overhead across all four options to answer whether PQC fits
within eUICC hardware and ES9+ bandwidth budgets.

\subsection{Two-Platform Testbed}\label{sec:meth-testbed}

No production eUICC yet supports ML-KEM or ML-DSA, so we combine two
platforms bridged by hardware-gap normalisation
(Section~\ref{sec:meth-normalization}).
\textbf{SM-DP+ server:} a modified osmo-smdpp~\citep{osmo_smdpp} linked
against liboqs~\citep{oqs2024} for ML-KEM-768 and ML-DSA-44, with
OQS-OpenSSL for PQ-TLS.
\textbf{Platform~A (sysmocom eUICC):} production eUICC via PC/SC over USB,
providing ground-truth classical baselines for all ECDSA and ECDH timings.
\textbf{Platform~B (Raspberry Pi Zero~2W):} ARM Cortex-A53 at 1\,GHz,
512\,MB RAM, running Raspberry Pi OS Lite (64-bit) with aarch64 liboqs;
ARMv8-A matches the SecurCore architecture of production eUICC silicon,
enabling meaningful cycle-count comparison.

\subsection{Metrics and Measurement Protocol}\label{sec:meth-metrics}

\textbf{Latency:} key-gen, sign, verify, encaps, decaps in milliseconds;
10{,}000 iterations after 100-iteration warm-up; mean, median, $\sigma$,
and 95\,\% CI reported; CPU frequency locked on Platform~B.
\textbf{Bandwidth:} per-certificate (CI root, sub-CA, DPauth, DPpb, EUICC,
EUM) and per-session byte counts across all four options.
\textbf{Memory:} peak heap during ML-DSA Sign and ML-KEM Decaps via valgrind
massif, compared to eUICC RAM budget.

\subsection{Hardware-Gap Normalisation}\label{sec:meth-normalization}

Classical operations are timed on both platforms; the per-operation scaling
factor is
\begin{equation}
\alpha_{i} = \frac{T_{\text{sysmocom}}(\text{op}_{i})}{T_{\text{Pi}}(\text{op}_{i})}
\end{equation}
PQC latency on the eUICC is projected as
$\hat{T}_{\text{eUICC}}(\text{PQC}) = \bar{\alpha} \times T_{\text{Pi}}(\text{PQC})$.
The validity of a pooled $\bar{\alpha}$ depends on how uniformly the two
platforms scale across operations.  We evaluate this via the coefficient of
variation $\text{CV} = \sigma(\alpha_i)/\mu(\alpha_i)$ at three levels:
globally across all six classical operations, and within asymmetric (ECDSA,
ECDH) and symmetric (AES, SHA) families.  The decision rule, fixed before
data collection, is:

\begin{itemize}[leftmargin=*, noitemsep, topsep=2pt]
  \item $\text{CV}_{\text{global}} \leq 0.10$: use a single pooled
        $\bar{\alpha}$; projection error $\lesssim$10\,\%.
  \item $\text{CV}_{\text{global}} > 0.10$ but per-category
        $\text{CV} \leq 0.10$: use separate $\bar{\alpha}_{\text{asym}}$
        and $\bar{\alpha}_{\text{sym}}$; PQC signature/KEM ops projected
        with $\bar{\alpha}_{\text{asym}}$, symmetric ops with
        $\bar{\alpha}_{\text{sym}}$.
  \item Per-category $\text{CV} > 0.10$: use per-operation $\alpha_i$;
        PQC ops with no classical analogue use
        $\bar{\alpha}_{\text{asym}}$ with an additional
        $\pm\text{CV}$ systematic uncertainty.
  \item $\text{CV}_{\text{global}} > 0.25$ with no category convergence:
        report Platform~B native timings only; linear normalisation is
        deemed unreliable.
\end{itemize}

\noindent Projection uncertainty is propagated as
$\sigma_{\hat{T}}^{2} = T_{\text{Pi}}^{2}\,\sigma_{\alpha}^{2} +
\alpha^{2}\,\sigma_{T_{\text{Pi}}}^{2}$, and reported with 95\,\% CIs.
ARM PMU cycle counts normalised to the 80--200\,MHz eUICC clock range
provide an independent cross-check.

% ===================================================================
%  SECTION 6 — FORMAL VERIFICATION
% ===================================================================
\section{Formal Verification}\label{sec:formal}

We model all four options as self-contained ProVerif~2.05
models~\citep{blanchet2016proverif} under a standard Dolev-Yao adversary
that controls the public channel.  The LPA is an honest passive relay.
Forward secrecy is tested via ProVerif's phase mechanism: long-term
signing keys are released in phase~1; the session key query is evaluated
in phase~0.

\subsection{Security Properties}\label{sec:fv-props}

Six properties derived from the SGP.22 security objectives and the HNDL
threat model:

\begin{enumerate}[leftmargin=*, noitemsep, topsep=2pt, label=\textbf{Q\arabic*.}]
  \item \textbf{Injective server auth:}  each eUICC acceptance of
        \texttt{serverSigned1} maps to a unique SM-DP+ send.
  \item \textbf{Injective client auth:}  each SM-DP+ acceptance of
        \texttt{euiccSigned1} maps to a unique eUICC send.
  \item \textbf{Session key secrecy:}  the KDF output \textit{ShS}
        is unreachable by the attacker.
  \item \textbf{Profile confidentiality:}  the plaintext eSIM profile
        is unreachable by the attacker.
  \item \textbf{Binding integrity:}  the eUICC installs a profile only
        if the authenticated SM-DP+ sent it (\texttt{bind\_body} MAC).
  \item \textbf{Forward secrecy:}  \textit{ShS} remains secret in
        phase~0 after long-term ECDSA/ML-DSA keys are revealed in
        phase~1.
\end{enumerate}

\subsection{ProVerif Encoding}\label{sec:fv-encoding}

ECDSA and ML-DSA-44 share the same EUF-CMA \texttt{reduc}; distinct
key types (\texttt{skey} vs.\ \texttt{mldsaskey}) distinguish them.
ECDH uses a symmetric Diffie-Hellman equation absent in option~(d).
ML-KEM-768 is split into two functions --- ciphertext and shared secret
--- so that \texttt{kem\_decaps} recovers the shared secret only from the
matching secret key, modelling IND-CCA2.

\begin{lstlisting}[caption={Cryptographic primitives}]
(* EUF-CMA sig -- ECDSA [a,b] or ML-DSA-44 [c,d] *)
fun pk(skey):pkey.
fun sign(bitstring,skey):bitstring.
reduc forall m:bitstring,k:skey;
  checksign(sign(m,k),pk(k))=m.
(* ECDH key agreement -- options (a)(b)(c) *)
equation forall x:dhskey,y:dhskey;
  dh(x,dhpk(y))=dh(y,dhpk(x)).
(* ML-KEM-768 IND-CCA2 -- option (d) only *)
fun kem_ct(kempkey,bitstring):bitstring.
fun kem_ss(kempkey,bitstring):bitstring.
reduc forall sk:kemskey,r:bitstring;
  kem_decaps(sk,kem_ct(kempk(sk),r))
           =kem_ss(kempk(sk),r).
\end{lstlisting}

In option~(d), \texttt{bind\_body} uses constructor
\texttt{mk\_bb\_kem(tid,\,ct,\,bpp)} --- carrying only the KEM
ciphertext and no SM-DP+ ephemeral DH key --- whereas options (a)--(c)
use \texttt{mk\_bb(tid,\,otPKdp,\,otPKeuicc,\,bpp)}, reflecting the
non-contributory structure of ML-KEM encapsulation.

\begin{lstlisting}[caption={Security queries (representative)}]
(* Q1: injective agreement on serverSigned1 *)
query ss1:bitstring;
  inj-event(eUICC_AuthenticatedServer(ss1))
  ==>inj-event(SMDP_SentServerSigned1(ss1)).
(* Q3: session key unreachable *)
query attacker(sessionKeyTest).
(* Q6: forward secrecy via phase *)
query attacker(sessionKeyTest) phase 0.
\end{lstlisting}

\subsection{Verification Results}\label{sec:fv-results}

All 24 queries (6 per option) return \textbf{true}.  Combined
verification time is under 3\,s (ProVerif~2.05, commodity x86-64).
Table~\ref{tab:fv-results} summarises results.  Q6 confirms
session-key independence from long-term keys for all options: ECDH
ephemeral keys in (a)--(c) and ML-KEM encapsulation randomness in (d)
both provide forward secrecy within the Dolev-Yao model.

\begin{table}[!htb]
\caption{ProVerif query results. \checkmark\,= verified \textbf{true}.}
\label{tab:fv-results}
\centering
\footnotesize
\setlength{\tabcolsep}{3pt}
\renewcommand{\arraystretch}{1.05}
\begin{tabularx}{\columnwidth}{@{} >{\raggedright\arraybackslash}X
    *{4}{>{\centering\arraybackslash}c} @{}}
\toprule
\textbf{Property} & \textbf{(a)} & \textbf{(b)} & \textbf{(c)} & \textbf{(d)} \\
\midrule
Q1 Inj.\ server auth      & \checkmark & \checkmark & \checkmark & \checkmark \\
Q2 Inj.\ client auth      & \checkmark & \checkmark & \checkmark & \checkmark \\
Q3 Session key secrecy    & \checkmark & \checkmark & \checkmark & \checkmark \\
Q4 Profile confidentiality& \checkmark & \checkmark & \checkmark & \checkmark \\
Q5 Binding integrity      & \checkmark & \checkmark & \checkmark & \checkmark \\
Q6 Forward secrecy        & \checkmark & \checkmark & \checkmark & \checkmark \\
\bottomrule
\end{tabularx}
\end{table}

% ===================================================================
%  SECTION 7 — PERFORMANCE RESULTS
% ===================================================================
\section{Performance Results}\label{sec:results}

\subsection{Experimental Setup}\label{sec:results-setup}

All experiments target the classical RSP channel (option~(b) transport
configuration): the SM-DP+ server presents a PQ-TLS endpoint and the
application-layer cryptography remains classical ECDSA/ECDH.
The testbed components are as follows.

\textbf{eUICC (Platform~A):} a production sysmocom eUICC (SGP.22 v2.3,
firmware~35.6.35) connected via a HID~OMNIKEY~3x21 PC/SC reader over USB.
The card reports 9{,}539\,B free volatile memory and 303\,KB free
non-volatile memory, confirming the constrained RAM profile that motivates
our normalisation methodology.  All per-APDU timings on this platform
include PC/SC host-stack overhead, which we calibrate and subtract
(Section~\ref{sec:results-cal}).

\textbf{SM-DP+ server:} osmo-smdpp~\citep{osmo_smdpp} compiled against
liboqs~\citep{oqs2024} for ML-DSA-44 certificate signing; fronted by an
nginx~1.26.3 instance linked against OpenSSL~3.6.0 for PQ-TLS (port~8443
for classical TLS~1.3, port~8444 for X25519MLKEM768\,+\,ML-DSA-44).  The
lpac LPA binary drives each provisioning session and instruments APDU/ES9+
timings via \texttt{BENCH} markers emitted on stderr.

\textbf{Measurement protocol:} 200 iterations per mode (classical then
PQ-TLS), each a complete profile download-and-delete cycle.  A 15-second
cooldown is inserted every 50 iterations.  Iterations with non-zero
\texttt{returncode} are excluded from latency statistics.

\subsection{PC/SC Overhead Calibration}\label{sec:results-cal}

We characterise the PC/SC stack latency as a linear function of payload
size by issuing UPDATE BINARY commands at five sizes
(0, 32, 64, 128, 255\,bytes; 500 iterations each).  Fitting yields

\begin{equation}
  T_{\text{PCSC}}(n) = 1{,}160\,\mu\text{s} + 62.4\,\mu\text{s}\,n,
  \quad R^{2} = 0.999,
\end{equation}

\noindent where $n$ is the APDU payload in bytes.  This model,
illustrated in Fig.~\ref{fig:pcsc-cal}, is subtracted from all
Platform~A timings to obtain net eUICC computation time.
The intercept (1.16\,ms) reflects fixed USB-HID round-trip cost; the
slope (62.4\,$\mu$s/byte) reflects the PC/SC framing and CCID
transport at USB~full-speed.

\begin{figure}[!htb]
\centering
\includegraphics[width=\columnwidth]{plots/pcsc_calibration.png}
\caption{PC/SC overhead model: latency vs.\ APDU payload size (500
  iterations per point). Linear fit $R^{2}=0.999$.}
\label{fig:pcsc-cal}
\end{figure}

\subsection{Session-Level Results: Classical vs.\ PQ-TLS}\label{sec:results-session}

Fig.~\ref{fig:session-wall} plots end-to-end session wall time for all
400 iterations (200 classical, 200 PQ-TLS).  Across both modes the mean
session duration is \textbf{25.20\,s} ($\sigma=0.67$\,s).
Table~\ref{tab:session-stats} quantifies the comparison.

\begin{table}[!htb]
\caption{End-to-end session wall time: classical vs.\ PQ-TLS
  (200 successful iterations each, mean\,$\pm\,1\sigma$).}
\label{tab:session-stats}
\centering
\footnotesize
\setlength{\tabcolsep}{5pt}
\renewcommand{\arraystretch}{1.05}
\begin{tabular}{@{}lrr@{}}
\toprule
\textbf{Mode} & \textbf{Mean (s)} & \textbf{$\sigma$ (s)} \\
\midrule
Classical TLS (option~a) & 25.20 & 0.66 \\
PQ-TLS transport (option~b) & 25.20 & 0.67 \\
\midrule
Difference & $<$0.01 & --- \\
\bottomrule
\end{tabular}
\end{table}

The two distributions are statistically indistinguishable: a Welch
$t$-test yields $p > 0.9$.  This result reflects the dominance of
on-card SIM processing over TLS transport overhead.

\begin{figure}[!htb]
\centering
\includegraphics[width=\columnwidth]{plots/session_wall_time.png}
\caption{Session wall time per iteration.  Blue: classical TLS (iters~1--200);
  purple: PQ-TLS (iters~201--400).  The dashed vertical line marks
  the phase transition.  Both modes sustain $\approx$25\,s/session.}
\label{fig:session-wall}
\end{figure}

\subsection{Operation-Level Breakdown}\label{sec:results-ops}

Fig.~\ref{fig:op-means} shows mean per-operation duration averaged over
all 400 iterations.  Two APDU commands account for the majority of
session time: \texttt{AuthenticateServer} (BF38, step~12,
3{,}121\,ms classical / 3{,}114\,ms PQ-TLS) and \texttt{PrepareDownload}
(BF21, step~20, 4{,}075\,ms / 4{,}112\,ms).  Together they consume
$\approx$28\,\% of the 25\,s budget.  The remaining time is distributed
across 18 \texttt{LoadBPP} segments (BF36), whose cumulative cost
($\sim$18\,s) reflects AES-128 decryption and integrity verification
performed entirely inside the eUICC.

\begin{figure}[!htb]
\centering
\includegraphics[width=\columnwidth]{plots/operation_means.png}
\caption{Mean APDU and ES9+ operation durations (all 400 iterations).
  On-card operations (APDU tags) dominate; ES9+ HTTP calls are
  $<$50\,ms each.}
\label{fig:op-means}
\end{figure}

The three ES9+ HTTP round-trips (\texttt{initiateAuthentication},
\texttt{authenticateClient}, \texttt{getBoundProfilePackage}) each take
14--37\,ms on average and are \textbf{identical} between classical and
PQ-TLS; their values are reported in Table~\ref{tab:es9p-times}.  TLS
transport cost, even with PQ key exchange, is fully absorbed within this
range.

\begin{table}[!htb]
\caption{ES9+ HTTP call latency: classical vs.\ PQ-TLS
  (mean over 200 iterations, ms).}
\label{tab:es9p-times}
\centering
\footnotesize
\setlength{\tabcolsep}{4pt}
\renewcommand{\arraystretch}{1.05}
\begin{tabular}{@{}lrr@{}}
\toprule
\textbf{ES9+ call} & \textbf{Classical (ms)} & \textbf{PQ-TLS (ms)} \\
\midrule
initiateAuthentication     & 14.7 & 12.2 \\
authenticateClient         & 36.5 & 33.6 \\
getBoundProfilePackage     & 26.4 & 23.8 \\
\bottomrule
\end{tabular}
\end{table}

\subsection{TLS Handshake Isolation}\label{sec:results-tls}

To isolate the bare TLS handshake overhead independent of APDU processing,
we run \texttt{openssl s\_client} in a tight loop (200 iterations each
mode) directly against the classical (port~8443) and PQ-TLS (port~8444)
nginx endpoints.  Fig.~\ref{fig:tls-handshake} presents the results.

\begin{figure}[!htb]
\centering
\includegraphics[width=\columnwidth]{plots/tls_handshake.png}
\caption{\textit{Left:} TLS~1.3 handshake latency on loopback (200
  iterations, mean~$\pm$~95\,\% CI).
  \textit{Right:} per-message byte breakdown of the handshake.
  Classical uses ECDSA P-256 certificate; PQ-TLS uses ML-DSA-44.}
\label{fig:tls-handshake}
\end{figure}

Handshake latency is 13.1\,ms (classical) vs.\ 12.3\,ms (PQ-TLS) on
loopback — within measurement noise.  The difference lies entirely in
message sizes:

\begin{itemize}[leftmargin=*, noitemsep, topsep=2pt]
  \item \textbf{Certificate:} 659\,B $\to$ 4{,}064\,B ($\times$6.2) — the
        ML-DSA-44 public key (1{,}312\,B) and added extension fields
        inflate the DER-encoded X.509 leaf.
  \item \textbf{CertificateVerify:} 79\,B $\to$ 2{,}428\,B ($\times$30.7)
        — the ML-DSA-44 signature replaces a 64-byte ECDSA-P256 value.
  \item \textbf{Total handshake:} $\approx$4.1\,KB $\to$ $\approx$9.7\,KB
        ($\times$2.4).
\end{itemize}

On a loopback interface transmitting at multi-gigabit rates, this $\sim$5.6\,KB
overhead is negligible in time.  On a narrowband link (e.g.\ 100\,kbps
NB-IoT), the same overhead adds $\approx$450\,ms per session
establishment — a cost that operators on constrained links must account
for in option~(b) deployment.

\subsection{Summary of Option~(b) Impact}\label{sec:results-summary}

Table~\ref{tab:option-b-summary} consolidates the measured option~(b)
overhead relative to the classical baseline.

\begin{table}[!htb]
\caption{Option~(b) overhead summary: PQ-TLS transport, classical
  app-layer.  Loopback measurements; 200 iterations each.}
\label{tab:option-b-summary}
\centering
\footnotesize
\setlength{\tabcolsep}{4pt}
\renewcommand{\arraystretch}{1.05}
\begin{tabular}{@{}lrrr@{}}
\toprule
\textbf{Metric} & \textbf{Classical} & \textbf{PQ-TLS} & \textbf{Ratio} \\
\midrule
Session wall time (s)      & 25.20 & 25.20 & 1.00$\times$ \\
TLS handshake latency (ms) & 13.1  & 12.3  & 0.94$\times$ \\
TLS total message size (KB)& 4.1   & 9.7   & 2.37$\times$ \\
Certificate size (B)       & 659   & 4{,}064 & 6.16$\times$ \\
CertificateVerify (B)      & 79    & 2{,}428 & 30.7$\times$ \\
\bottomrule
\end{tabular}
\end{table}

The central finding is that \textbf{option~(b) imposes zero measurable
latency overhead} on end-to-end RSP sessions: the $\sim$5.6\,KB increase in
TLS handshake bytes is absorbed entirely within the dominant on-card SIM
processing time ($\sim$25\,s).  The real deployment cost of option~(b) is
therefore bandwidth-, not computation-bound, and is relevant only for
low-throughput ES9+ links.

% ===================================================================
%  SECTION 7 — DISCUSSION AND CONCLUSION
% ===================================================================
\section{Discussion and Conclusion}\label{sec:conclusion}

The authentication phase (steps 1--16) concentrates the heaviest PQC
overhead: four ML-DSA sign/verify operations and two full CI certificate chain
transmissions under options~(c) and~(d), while the download phase adds ML-KEM
encaps/decaps only in option~(d) at substantially lower computation cost.
Option~(b) is immediately deployable as a purely server-side change that
provides HNDL protection for the ES9+ transport channel without any eUICC
modification.  Option~(c) addresses the higher-priority authentication-forgery
threat once new ML-DSA-capable eUICC silicon becomes available, deferring the
key-exchange migration.  Option~(d) achieves full application-layer quantum
resistance but requires eUICC support for both ML-DSA and ML-KEM, CI root
re-issuance, and SGP.22 specification updates for PQC algorithm identifiers
in \texttt{euiccInfo1/2} and fragmentation support for enlarged certificate
chains.  The NSA CNSA~2.0 mandate (2035)~\citep{nsa_cnsa2} implies that
devices manufactured today with typical 5-year lifetimes will still be active
past the deadline; operators should treat option~(b) as an immediate baseline
and plan eUICC refresh cycles around options~(c) and~(d) in alignment with
GSMA PQ.04~\citep{gsma_pq04} and NIST IR~8547~\citep{nist_ir8547}.

Section~\ref{sec:formal} demonstrates through ProVerif that all six
security properties — including session-key forward secrecy — hold across
all four options; ML-KEM encapsulation in option~(d) provides the same
Dolev-Yao-level guarantees as ephemeral ECDH in options (a)--(c) while
eliminating the classical key-agreement vulnerability.
Future work includes benchmarking on PQC-capable eUICC silicon once
available and evaluation of hybrid certificate formats for a gradual CI
trust-hierarchy transition.

\section*{Conflict of interest}
The authors declare no competing interests.

\bibliographystyle{elsarticle-num}
\bibliography{sample}

\end{document}
